\section{Connectedness and potential pro-modularity}
\label{sec:connectedness}

We now connect a modular seed to the target. The field
construction of Section~\ref{sec:fields} makes the reducible
locus small after recording local lines at selected dyadic
primes. The unrestricted propagation theorem of
Section~\ref{sec:propagation} can then cross successive
component intersections.

\begin{theorem}\label{thm:potential-promodularity}
Let \(r\) satisfy the hypotheses of Theorem~\ref{thm:main} and be not
virtually solvable. Then \(r\) is pro-modular over
some solvable totally real Galois extension of
\(\mathbb Q\), of even degree and split completely
at \(2\).

More generally, fix a complete integral framed deformation family over
\(\mathbb Q\), with determinant \(\chi=\det r\), containing the
coefficient point defined by \(r\), and with ramification confined
to a finite set. The closed trace image of this family is pro-modular
after restriction to one common extension of the stated kind, at
one fixed tame level.
\end{theorem}

To prove the family assertion, we will compare a curve in a Hecke
component with a curve in an unrestricted deformation component
containing the target family. The first curve requires a cuspidal
lift of the target's semisimple residue. We choose that lift before
preparing the fields, so that its tame primes belong to the fixed
allowed set used in the field estimates.

\subsection{A modular residual seed}

\begin{lemma}\label{lem:modular-seed}
Every semisimple two-dimensional residual
representation arising from the target has a
classical cuspidal lift over \(\mathbb Q\), after
finite coefficient extension and allowing a
finite enlargement of the tame level.
\end{lemma}
\begin{proof}
For absolutely irreducible residue this is
Serre's modularity theorem, including
characteristic two
\cite[Theorem~0.2]{Kisin2005Survey};
the results are due to
Khare--Wintenberger and Kisin
\cite{KhareWintenberger2009,KhareWintenberger2009II,Kisin2009}.
The characteristic-two parity condition is
automatic.

For reducible residue, the existence of a cuspidal lift, including
in characteristic two, is already proved by Billerey--Menares
\cite[Theorem~2.1]{BillereyMenares2016}. The geometric argument below
uses the same Eisenstein difference as their
\cite[Theorem~2.3]{BillereyMenares2016}; we retain the details of
integrality, cuspidality, and lifting needed here.

Write the reducible residue as
\(\bar\chi_1\oplus\bar\chi_2\) and put
\(\alpha=\bar\chi_1/\bar\chi_2\). After a finite-order twist
it suffices to lift the eigensystem of
\(1\oplus\alpha\).
The character \(\alpha\) has odd order,
is unramified at \(2\), and has an even
Dirichlet lift of odd conductor.
Choose an even weight \(k'\ge4\) and a normalized
Eisenstein series \(f\) with this character pair
and \(a_1(f)=1\). Take an odd level \(N\ge5\)
divisible by the conductor. Let \(\ell\) be a
prime congruent to one modulo \(N2^a\), where
\(a\) will be sufficiently large. Put
\[
 g(\tau)=f(\tau)-f(\ell\tau).
\]
Regard this as a form of level \(\Gamma_1(N\ell)\).
Its positive Fourier coefficients at infinity
are \(2\)-integral, and its constant term there
is zero. The tame modular curve has good
reduction, with geometrically integral special
fiber \cite[Corollary~10.9.5(2) and Table~10.9.6]{KatzMazur1985}.
The \(q\)-expansion principle
\cite[Theorem~1.6.1 and Corollary~1.6.2]{Katz1973} therefore
extends \(g\) integrally: a vertical pole would
be detected at that component's cusp expansion.

We verify cuspidality modulo \(\varpi\) at every
cusp. For a cusp matrix \(\gamma\in\SL_2(\mathbb Z)\),
the primitive-column decomposition is
\[
 \begin{pmatrix}\ell&0\\0&1\end{pmatrix}\gamma
   =\gamma'
     \begin{pmatrix}e&t\\0&\ell/e\end{pmatrix},
 \qquad e\in\{1,\ell\}.
\]
Since \(\ell\equiv1\pmod N\), one can choose
\(\gamma'\equiv\gamma\pmod N\).
The two constant terms in \(g|_{k'}\gamma\)
therefore differ by the factor
\[
 1-(\ell/e)^{-k'}.
\]
There are only finitely many constant terms of
\(f\) at level \(N\). Choose \(a\) larger than
their denominator valuations, with an additional
power sufficient for reduction modulo
\(\varpi\). The displayed factor makes every
constant term of \(g\) vanish in the residue
field. Its coefficient of \(q\) is one, so the
reduction is nonzero. Away from \(N\ell\), it
is an eigenvector with the Eisenstein residual
eigensystem \(1\oplus\alpha\).

This mod-\(\varpi\) cusp form lifts to the
characteristic-zero cuspidal space in the same
weight. Indeed on the tame modular curve,
Kodaira--Spencer gives
\(\Omega^1\simeq\omega^2(-\mathrm{cusps})\).
For this formula and positivity of \(\omega\), use
\cite[Section~1.5 and Appendix~A1.5.4]{Katz1973}:
after adjoining the required roots of unity, a symplectic
full-level component maps to the \(\Gamma_1\)-curve
by a finite cover of odd degree. The logarithmic
Kodaira--Spencer isomorphism descends along this tame
cover, and positivity descends along the finite map.
On the special fiber, Serre duality
\cite[Lemma~53.4.1 and Remark~53.4.3, Tags~0BS2 and~0BS4]{StacksProject}
identifies the obstruction
\(H^1(\omega^{k'}(-\mathrm{cusps}))\) with
the dual of \(H^0(\omega^{2-k'})\), which is
zero for \(k'\ge4\). The corresponding integral \(H^1\)-module
is finite by proper coherent cohomology
\cite[Lemma~30.19.2, Tag~02O6]{StacksProject}.
The multiplication-by-\(\varpi\) exact sequence and the special-fiber
vanishing make multiplication by \(\varpi\) surjective on this module.
Nakayama's lemma makes it zero, so reduction of cuspidal sections is
surjective.
The finite integral cuspidal Hecke module has
a characteristic-zero eigenpacket specializing
to this residual eigensystem: its Hecke algebra
is finite and torsion-free over \(\OO\), and
the residual maximal ideal lies above a
horizontal minimal prime. Finally undo the
finite-order twist. The construction also
covers \(\alpha=1\).
\end{proof}

\subsection{Field choices}

Fix the target's semisimple residue \(\bar r\) and determinant
\(\chi\). Choose a cuspidal lift of this residue by
Lemma~\ref{lem:modular-seed}, and fix a finite allowed set over
\(\mathbb Q\) containing the bad primes of the target, the prescribed
family when present, and this lift. This choice precedes the field construction below;
it ensures that the bound for the number of allowed primes also
covers the modular seed. In the reducible case write
\(\bar r^{\mathrm{ss}}=\bar\chi_1\oplus\bar\chi_2\) and
\(\alpha=\bar\chi_1/\bar\chi_2\). Use the fields
\[
 \mathbb Q\subset A\subset F
\]
of Proposition~\ref{prop:field-preparation}, with
\(d_A=[A:\mathbb Q]\), \(d=[F:\mathbb Q]\).
There are no marks when \(\bar r\) is absolutely
irreducible or \(\alpha|_{G_{\mathbb Q_2}}\ne1\).
Otherwise choose the distinguished place \(v_0\) of \(A\)
and mark its \(m=[F:A]\) lifts to \(F\).
The numbers \(u_A,u_F\) of allowed non-dyadic primes
are uniformly bounded as these degrees grow.

Choose the degrees sufficiently large that
\begin{equation}\label{eq:degree-choices}
 d>u_F+4,\qquad
 m>u_F+4\quad\text{if there are marks},\qquad
 2d-2m-4-u_F>d.
\end{equation}
In the unmarked case set \(m=0\).
Since \(m/d=1/d_A\) in the marked case and \(m/d=0\)
otherwise, the further fixed lower bounds on \(d_A\)
below can be imposed simultaneously.

\subsection{Charts with transverse local lines}

Over the characteristic-two fixed-determinant pseudo-base,
take the finite-type framed representation scheme and
adjoin a local invariant projective line at every mark.
Call the resulting scheme \(Y\).
In its closed residual fiber, call a point \emph{transverse}
if none of its marked lines is globally invariant.
The condition is geometric: it must hold after extending
its residue field to an algebraic closure.
There is no additional condition when there are no marks.
The transverse locus is open in the residual fiber,
since the incidence of a global invariant line equal
to a specified marked line is proper.

At a transverse point with marks, the local residual
representation at each mark is scalar after untwisting.
Indeed there is a global invariant line, and the
different marked invariant line splits its restriction.
The two local semisimple characters coincide.
This simple observation is responsible for the strict
Selmer condition in the next lemma.

\begin{lemma}\label{lem:good-connectedness}
Let \(F\), the marked set, and the fixed-determinant representation
scheme \(Y\) be as above, with the degree inequalities
\eqref{eq:degree-choices}. Let \(y\) be a transverse closed residual
point of \(Y\), with finite residue field, and let \(R'_y\) be
the completion of \(Y\) at \(y\). Every component of
\(\Spec R'_y\) has dimension at least \(2d-m\).
Its components can be joined by a chain whose successive
intersections have dimension at least \(2d-m-1\).
The globally reducible locus has dimension strictly
less than \(2d-m-1\).
\end{lemma}
\begin{proof}
Retain the real local problems and the \(m\) marked
dyadic line problems as the local base. Reducing this
base changes no spectrum. By
Lemma~\ref{lem:marked-local-base}, it is a complete
local domain of dimension
\[
 2d+5m.
\]
The relative framed presentation of
Proposition~\ref{prop:relative-presentation} has
generators minus relations at least
\[
 3-6m-h^0((\ad^0)^\vee(1))\ge-6m.
\]
The last inequality uses the three-dimensionality of
\((\ad^0)^\vee(1)\). A quotient of a complete local
domain with this presentation has all components of
dimension at least \(2d-m\). Grothendieck's
connectedness theorem gives connectedness in dimension
at least \(2d-m-1\): in the form used here, cutting a
complete local domain of dimension \(N\), after adding
\(a\) formal variables, by \(b\) equations gives
connectedness in dimension at least \(N+a-b-1\).
This is the component-chain formulation of the
connectedness theorem \cite[Expos\'e~XIII, Th\'eor\`eme~2.1]{SGA2}.

We bound the reducible locus separately.
It is empty for absolutely irreducible residue.
Otherwise add a global invariant line by proper
incidence. At a closed point of that incidence,
choose an adapted basis and first hold this line
fixed. Its movement costs one parameter.
The diagonal global characters with fixed product
vary in dimension at most one, by abelian
Leopoldt over \(F\); see
Lemma~\ref{lem:pseudo-bad-loci}.
Freeze those character parameters.

At a marked place the closed local matrices are
scalar and the marked line is transverse to the
chosen global line. An infinitesimal upper extension,
with diagonal characters frozen, must therefore
vanish locally: a nonzero upper nilpotent matrix
cannot preserve a second transverse line.
Thus its cohomology class is strict at the marks.
Moving a marked line costs at most one further
tangent parameter, giving \(m\) in total.
There is at most one upper coboundary parameter.
The strict Selmer bounds in \eqref{eq:strict-selmer}
therefore give the following dimensions:
\[
 \begin{cases}
 d+u_F+3,&\alpha|_{G_{\mathbb Q_2}}\ne1,\\
 2d-2m+u_F+3,&\alpha|_{G_{\mathbb Q_2}}=1.
 \end{cases}
\]
Here the three remaining parameters account for
the diagonal base, the global line, and the
possible coboundary. The relative tangent estimate
bounds the number of generators over the diagonal
base, so it is an upper bound for the whole formal
reducible chart, not only for its tangent space at
a fixed diagonal character. Properness of the
global-line incidence transfers the bound to
its image. Finite residue extensions do not
affect these dimensions.

The first two inequalities in
\eqref{eq:degree-choices} make both bounds smaller
than \(2d-m-1\). This proves the lemma.
\end{proof}

\begin{lemma}\label{lem:chart-propagation}
At a transverse closed residual point, either every
component of \(Y\) through that point has potentially
pro-modular pseudo-image, or none does. It suffices
to begin with one component whose pseudo-image
contains a potentially pro-modular,
non-virtually-solvable curve having finite local
images at the allowed non-dyadic places.
\end{lemma}
\begin{proof}
Work first in the completed chart.
By Lemma~\ref{lem:good-connectedness}, two successive
components in a connecting chain intersect in
dimension at least \(2d-m-1\).
Choose an irreducible component of maximal dimension
in that intersection. Its generic global representation
is absolutely irreducible, because the reducible locus
has smaller dimension.
Pass first to the completed trace-and-character image,
retaining the characters on the marked lines.
Forgetting the global frame and the line choices loses
at most \(3+m\) dimensions, by
Lemma~\ref{lem:dimension-comparison}.
Proposition~\ref{prop:rep-space} makes this image finite
over its trace image, so forgetting the recorded characters
causes no further dimension loss.
Thus the pseudo-image of the chosen component has dimension at least
\[
 2d-2m-4.
\]
At each allowed non-dyadic prime impose that the
trace of a chosen Frobenius lift equal its residual
value. These \(u_F\) equations leave a closed locus
of dimension at least
\(2d-2m-4-u_F>d\).
The bad pseudo-locus bounds of
Lemma~\ref{lem:pseudo-bad-loci} and curve avoidance
therefore produce a non-virtually-solvable curve
in it. Lemma~\ref{lem:inertia-basechange} gives
finite local images at the non-dyadic allowed
places.

If the preceding component is potentially
pro-modular, this curve is potentially pro-modular.
Theorem~\ref{thm:localized-propagation} makes the
next component potentially pro-modular.
A starting curve as in the statement first makes
its containing component potentially pro-modular
by that same theorem. Iterating crosses the
entire chain. Only finitely many extensions and
levels occur. Take the compositum of their Galois
closures and a common deeper level, using
Lemma~\ref{lem:forward-basechange}.

Completion is faithfully flat. Every component
of \(Y\) through the closed point is the
contraction of a minimal prime of the completed
chart. The component pseudo-images are the
corresponding trace contractions. The conclusion
therefore descends from the chart to \(Y\).
\end{proof}

The local chart statement must next be joined across
different residual extension classes. The transverse
residual condition was chosen so that this joining
does not require changing the semisimple residue.

\begin{lemma}\label{lem:transverse-fiber}
If nonempty, the transverse closed residual fiber
of \(Y\) is geometrically connected. Consequently,
if one transverse chart has potentially pro-modular
component images, every transverse chart does.
\end{lemma}
\begin{proof}
Work over an algebraic closure of the residue field.
For scalar residue with marks, a transverse residual
representation is a nonzero additive extension after
a scalar twist. It has a unique global invariant line.
Its cocycle is zero at the marks. With the global
line fixed, these cocycles form a vector space with
the origin removed. The marked complementary lines
form products of affine line charts. Conjugating
the basis gives a connected parameter space whose
image is the entire transverse fiber.

For distinct residual characters, use both upper
triangular orientations. In either orientation,
with its global line fixed, the extension cocycles
form a vector space. In the marked case they
vanish locally at the marks. The allowed choices
of marked lines form a nonempty open subset of
the product with projective lines, hence an
irreducible parameter space. Varying the basis
preserves connectedness. The two orientations
meet at the split representation; when there
are marks, choose lines avoiding both global
invariant lines. Without marks there is no
transversality condition, and the same two
families meet at that split representation.
For absolutely irreducible residue, the
residual fiber is a conjugacy orbit.
These parameter spaces are finite-dimensional
algebraic spaces of cocycles: this follows either
from the residual representation scheme or from
the finite-dimensional residual \(H^1\).
This proves geometric connectedness.

There are finitely many components of the noetherian
scheme \(Y\). Their nonempty intersections with the
transverse residual fiber form a finite closed cover.
Make a graph with these components as vertices, joining
two vertices when they meet in that fiber. The graph
is connected because the fiber is connected.
Choose a spanning tree and, for each edge, a closed point
of the corresponding intersection in the transverse locus,
after a finite coefficient extension. At that point
Lemma~\ref{lem:chart-propagation} gives the same potential
pro-modularity status to the two components.

Only finitely many such points and completed-chart
component chains are needed. Their field extensions
have a common Galois compositum, and forward base change gives
one common deeper tame level. Thus the conclusion
holds simultaneously for every component meeting the
transverse fiber, and hence for every transverse chart.
Every component through a transverse point is accounted
for by completion, as in the preceding lemma. The
finite coefficient extensions are harmless because the
integral forms and their Hecke support commute with them.
\end{proof}

\subsection{From the two curves to the target family}

The proof first constructs large seed and target loci over \(A\),
then chooses a characteristic-two curve in each. Appropriate lattices
put both curves in transverse charts over \(F\). Connectedness transfers
potential pro-modularity from the seed curve to the target curve;
unrestricted propagation then reaches the entire containing target locus.

\begin{proof}[Proof of Theorem~\ref{thm:potential-promodularity}]
\emph{The seed and target loci over \(A\).}
For the seed and fields already chosen, solvable base change and Jacquet--Langlands
give a cuspidal residual packet over \(A\).
It remains cuspidal by regularity, as in
Lemma~\ref{lem:forward-basechange}.

Its central character need agree with
\(\psi=\chi\eps\) only modulo \(\varpi\).
At a sufficiently deep tame level their
transformations coincide on reduction.
The exact integral function model of
Proposition~\ref{prop:completed-forms} therefore
gives residual occurrence also for the
specified central character \(\psi\).
Proposition~\ref{prop:hecke-dimension} supplies
a pro-modular component \(C_{\mathrm{seed}}\)
in the fixed-determinant pseudo-base over \(A\)
with
\[
 \dim C_{\mathrm{seed}}\ge1+2d_A.
\]

Independently choose an unrestricted global
fixed-determinant framed component through
\(r|_{G_A}\), or containing the restricted
framed support of the prescribed family.
Such a component exists by taking a minimal
prime below that support. Its generic
representation is absolutely irreducible
because it contains the absolutely irreducible
target point.
The relative presentation over the real
local rings gives dimension at least
\[
 1+2d_A+3-h^0((\ad^0)^\vee(1))
\]
before forgetting the frame. The loss under
trace contraction is at most three.
Its pseudo-image \(C_{\mathrm{target}}\)
therefore has dimension at least
\(2d_A-2\). This large unrestricted component
contains the trace image of the initially
prescribed family; no lower dimension
bound on that smaller family is required.

\smallskip
\noindent\emph{Choosing curves in the two loci.}
In each of \(C_{\mathrm{seed}}\) and
\(C_{\mathrm{target}}\), pass to the special
fiber and prescribe residual Frobenius trace
at every bad non-dyadic place. If there are
marks, prescribe in addition that the
entire local pseudorepresentation at
\(v_0\) is its constant residual one.
This last condition costs at most the
number of topological generators of one
fixed local pseudo-ring for \(\mathbb Q_2\),
a constant independent of \(d_A\).
All the other costs are bounded by
\(1+u_A\), with a fixed additional constant
for the target dimension estimate.
Taking \(d_A\) sufficiently large leaves
dimension greater than \(d_A\).
By Lemma~\ref{lem:pseudo-bad-loci}, each locus
contains a non-virtually-solvable curve.
The Frobenius constraints make its bad
local images finite.

\smallskip
\noindent\emph{Entering the transverse fiber over \(F\).}
We explain why both curves enter transverse
charts over \(F\).
Without marks, any stable lattice suffices.
With marks, the local pseudo at \(v_0\) is
constant and scalar after untwisting,
so the curve-field representation has an
invariant line there after a finite
coefficient extension.
Choose a nonzero vector \(v\) on this line
and let
\[
 \mathcal V=\sum_{g\in G_A} k'[[t]]\,\rho(g)v.
\]
Compactness bounds this module inside a fixed
lattice. Global irreducibility makes it
full rank, and it is stable by construction.
The vector \(v\) is primitive in
\(\mathcal V\): if it belonged to
\(t\mathcal V\), all its translates would,
and Nakayama would force
\(\mathcal V=0\).
Its reduction lies in no global invariant
line, since its orbit spans
\(\mathcal V/t\mathcal V\).

Restrict the lattice to \(G_F\). No new global invariant
line appears in the residue. For distinct residual
characters, a nonsplit extension remains nonsplit by
the restriction injectivity in
Proposition~\ref{prop:field-preparation}; a split
representation retains its two distinct character lines.
For scalar semisimplification, the spanning property
above rules out a split scalar residue, and its nonzero
additive extension remains nonzero by the same injectivity.
These descriptions hold after extending the residue field.
Since \(F/A\) is Galois, \(G_A\) permutes the globally
invariant lines of the restricted residue. Transporting
the chosen local line to every prime over \(v_0\) therefore
preserves its avoidance of all these lines. Hence every
transported marked line is transverse.
Both curves pass through transverse points
of the same residual scheme \(Y\).

\smallskip
\noindent\emph{Returning from the target curve to its containing locus.}
The seed curve remains potentially
pro-modular after restriction to \(F\).
Its finite bad local images allow
Theorem~\ref{thm:localized-propagation} to
start the propagation in its transverse
chart. Lemma~\ref{lem:transverse-fiber}
then gives potentially pro-modular images
at every transverse chart.
In particular the restricted target curve is potentially
pro-modular. By definition, the original target curve
over \(A\) is therefore potentially pro-modular as well.
Apply Theorem~\ref{thm:localized-propagation}
over \(A\) to this original curve and its containing
locus \(C_{\mathrm{target}}\).
The whole locus, and therefore the target
and the prescribed family trace image,
becomes pro-modular after a common allowed
extension.

The finite component graph in
Lemma~\ref{lem:transverse-fiber} and the finite chains
inside its selected charts require only finitely many
propagation steps. Replace their fields by one Galois
compositum and their levels by one deeper tame level,
using forward base change. Total reality, complete
splitting at \(2\), and solvability over \(\mathbb Q\)
are preserved.
This proves both assertions.
\end{proof}

The distinction between the two spaces in this proof
is important. Marks and transverse lines produce
connectedness in the global special fiber.
The propagation theorem subsequently applies to
unrestricted dyadic deformation loci.
No ordinary or potentially crystalline restriction
has been imposed on the target. The next section
extracts its classical vectors.
